\subsection{Short-first-pair ordinary eighth moments for odd terminals of minimum two}\label{proof:Q5-C095}

\paragraph{Q95.0 Structures and bounds.}
The receiving root has an outward unit arm, all connector edges point
from the terminal to this root, and \(L\ge\allowbreak 1\). For genuine equal-pair
products \(E_t,\allowbreak E_r\), put \(H=\allowbreak E_rpT^L(E_tF)\). We prove the following
three rows, for every finite arm list and every positive connector length:
\fitdisplay{
\begin{array}{c|c|c}
\text{distinct residual arms}&F&\text{upper bound for }\E H^8\\\hline
(a,4,c),\ a=2,3,\ c\ge5&\bar f_ah_4f_c&e^{-d}\\
(2,3,4)&\bar f_2f_3h_4&e^{-3d/32}\\
2<3<c_1<b_2<c_2<\cdots<b_k<c_k,\ k\ge2
&\bar f_2f_3h_{c_1}\prod_{i=2}^kh_{b_i}f_{c_i}
&e^{-[1+(128/5)(k-2)]d}
\end{array}
}{}
The last row with \(c_1=\allowbreak 4\) has the stronger bound
\(e^{-[5+(128/5)(k-2)]d}\).
The \(\bar f_a\)-arm is fully inward, each \(f_n\)-arm fully outward,
and each \(h_n\)-arm has its first edge inward and all later edges
outward. In the last two rows the first pair is \(f_3h_{c_1}\).
Directions are fixed by the integer arm list and are the same in
the low and high parameter domains.

\paragraph{Q95.1 The finite inputs.}
Use \(A=\allowbreak J+K\), \(T=\allowbreak J-K=\allowbreak A^*\), \(g=\allowbreak K\one\), \(p=\allowbreak 1+g\), \(q=\allowbreak 1-g\),
\(f_n=\allowbreak A^n\one\), \(\bar f_n=\allowbreak T^n\one\), \(h_n=\allowbreak Tf_{n-1}\), and
\(\delta_i=\allowbreak \|K^i\one\|^2\), \(d=\allowbreak \delta_1\), as in
Appendix~\ref{proof:Q5-C064}. In the fixed measure,
\begin{equation}\tag{Q95.1.1}\label{eq:Q95:1.1}
\begin{gathered}
0\le d\le1/3,\quad d^2\le\delta_2\le\theta d,\quad
\delta_{i+1}\le\theta\delta_i,\quad\|K\|^2\le\theta,\\
|Kv(x)|\le\|v\|_2,\quad\theta=13/32,\quad\kappa=19/32,\\
0\le f_n,\bar f_n,h_n\le2,\quad\mu_n=\E f_n=\E\bar f_n\in[0,1],\\
\quad\E K^{2j}\one=(-1)^j\delta_j,\quad\E K^{2j+1}\one=0.
\end{gathered}
\end{equation}
Each genuine pair satisfies \(f_\ell+h_\ell=\allowbreak 2\mu_{\ell-1}\) and
\(0\le\allowbreak  f_\ell h_\ell\le\allowbreak 1\), and can be restored by positivity.
After removal of zero-mass atoms, \(p(x),\allowbreak q(x)\ge\allowbreak \mu(x)>\allowbreak 0\).
The all-equal labelling makes every tree message positive.
For \(d=\allowbreak 0\), \(g=\allowbreak 0\) and leaf integration gives one, including when
\(K\ne0\); divisions below are restricted to \(d>\allowbreak 0\).
Set \(\tau=\allowbreak d-\delta_2\), \(\beta_1=\allowbreak \beta_2=\allowbreak 0\), and
\begin{equation}\tag{Q95.1.2}\label{eq:Q95:1.2}
\begin{gathered}
\beta_n=\sum_{j=1}^{\lfloor(n-1)/2\rfloor}(-1)^{j-1}\delta_j
\quad(n\ge3),\quad\tau\le\beta_n\le d,\quad\tau\ge\kappa d,
\quad\beta_3=\beta_4=d,\\
S_n=\sum_{j=0}^nK^j\one=f_n+\sum_{i=3}^n\beta_i f_{n-i}.
\end{gathered}
\end{equation}
Under \(K'=\allowbreak -K\), \(g'=\allowbreak -g\), the convolution for \(S'_n\) uses
\(\bar f_n\). With \(G_m(z)=\allowbreak \sum_{j=\allowbreak 0}^{m-1}z^j\),
\(\|G_m(K)\|\le\allowbreak 1\) for \(m=\allowbreak 0,\allowbreak 1\) or \(m\ge\allowbreak 3\), whereas
\(\|G_2\|\le\allowbreak \sqrt{1+\theta}<\allowbreak 6/5\).
The full-mean exponential inequality of Appendix~\ref{proof:Q5-C081} is
\begin{equation}\tag{Q95.1.3}\label{eq:Q95:1.3}
\E e^{8K'V}\le e^{8\E(K'V)+17\|V\|_2^2}
\end{equation}
for each real \(V\).
For a fixed row, write \(k(y)=\allowbreak K(x,\allowbreak y)\). Skewness gives
\fitdisplay{
r(x)=Kg(x)=\frac12\E_{y,z}[(k(y)-k(z))K(y,z)].
}{}
Since \(|u-v|\le\allowbreak 1-uv\) for \(u,\allowbreak v\in[-1,\allowbreak 1]\),
\(|r(x)|\le\allowbreak (1-g(x)^2)/2=\allowbreak pq/2\); hence \(f_2=\allowbreak p+r\ge\allowbreak  p^2/2\).
For every \(5\le\allowbreak  s<\allowbreak t\), the all-index path mean
\(\mu_n\le\allowbreak 1-83d/32\) for \(n\ge\allowbreak 4\) and the positive finite-difference
bound \(4\delta_2/(1-\theta^2)\le\allowbreak 1664d/855\) give pair-sum loss
\([83/16-1664/855]d=\allowbreak [16/5+113/2736]d\).
AM--GM therefore gives \(h_sf_t\le\allowbreak  e^{R_{st}-16d/5}\), where
\(R_{st}=\allowbreak \sum_{r=\allowbreak s+1}^tK^r\one\); see Appendix~\ref{proof:Q5-C090}.

\paragraph{Q95.2 Finite propagation retaining the inward arm.}
Here \(a=\allowbreak 2\) or \(3\). Direct expansion gives
\(\bar f_2=\allowbreak q+Kg\), \(\E\bar f_2=\allowbreak 1-d\), and
\(\sigma(\bar f_2)^2=\allowbreak d+\delta_2-d^2\le\allowbreak 36d/25\).
The orthogonal odd/even decomposition
\(\bar f_3-\E\bar f_3=\allowbreak -(g+K^3\one)+(Kg+d)\) gives
\(\sigma(\bar f_3)^2=\allowbreak d-\delta_2+\delta_3-d^2\le\allowbreak  d\).
Thus \(\sigma(\bar f_a)\le\allowbreak 6\sqrt d/5\).
Suppose $0\le F\le\bar f_a e^{R-\Gamma}$ and $|R|\le1/2$.
Set $\mu=\E\bar f_a$ and $D_i=\E B^iR$, with $B=K'=-K$ and $u=B\one=g'$.
Apply Lemma~\ref{lem:oc:anchored-eighth} with
$s=\sigma(\bar f_a)/\sqrt d$ and $\varrho=\|R\|/\sqrt d$.
The first step and its finite continuation are
\begin{equation}\tag{Q95.2.1}\label{eq:Q95:2.1}
 T(\bar f_ae^R)\le\bar f_{a+1}+\mu BR+\delta,\qquad
 \delta=\mu D_0+2\sigma(\bar f_a)\|R\|+\tfrac{12}{5}\|R\|^2,
\end{equation}
\begin{equation}\tag{Q95.2.2}\label{eq:Q95:2.2}
 0<e^\Gamma T^LF\le Y_L=\bar f_{a+L}+\mu B^LR
 +\delta\bar f_{L-1}+\mu\sum_{i=1}^{L-1}D_i\bar f_{L-1-i}.
\end{equation}
For an identically zero input replace the initial strict inequality by
a weak one. The moment bound follows by continuity. Use the common
parameters, here written explicitly with $n=a+L$:
\[
\begin{gathered}
 V=G_{n-1}(B)u+\mu B^{L-1}R,\quad N=BV,\quad
 \E N=-\beta_{n+1}+\mu D_L,\\
 c_1=\delta-\beta_n,\qquad
 c_L=2\delta_++2\mu\sum_{i=1}^{L-1}(D_i)_+-\beta_n\quad(L\ge2).
\end{gathered}
\]
Thus
\begin{equation}\tag{Q95.2.3}\label{eq:Q95:2.3}
 \log\E[(pT^LF)^8]\le-8\Gamma+8c_L+8\E N+17\|V\|^2.
\end{equation}
The actual messages are positive. At $d=0$ use leaf integration rather
than the displayed ratios; for a zero input the logarithm is $-\infty$.

\paragraph{Q95.3 Triples \((a,\allowbreak 4,\allowbreak c)\), \(a=\allowbreak 2,\allowbreak 3\), \(c\ge\allowbreak 5\).}
Take \(F=\allowbreak \bar f_ah_4f_c\). Retaining the negative \(f_2\) term gives
\fitdisplay{
h_4+f_c\le2+R-4d+\delta_2(1+p)-\tau f_2,\qquad
R=\sum_{r=5}^cK^r\one,\quad f_2\ge p^2/2.
}{}
Completing the square bounds the excess by
\(\delta_2+\delta_2^2/(2\tau)\).
For \(t=\allowbreak \delta_2/d\le\allowbreak \theta\), the function
\(t+t^2/[2(1-t)]\) increases on \([0,\allowbreak \theta]\), so this is at most
\(663d/1216\). The pair loss is at least
\(4201d/1216>\allowbreak 69d/20\), with difference \(29d/6080\). Hence
\(F\le\allowbreak \bar f_ae^{R-\Gamma}\), \(\Gamma=\allowbreak 69d/20\).
The interval is consecutive:
\(R=\allowbreak K^4G_{c-4}(K)g\). Therefore
\fitdisplay{
\|R\|\le(6/5)\theta^2\sqrt d=507\sqrt d/2560<\sqrt d/5,
\qquad |R|<\sqrt d/3\le1/3.
}{}
The Taylor bound holds on the whole domain. Alternating interval
moments give \(D_0\le\allowbreak 0\), \(D_1\le\allowbreak \delta_3<\allowbreak d/6\), and
\(D_i\le\allowbreak \delta_4<\allowbreak 7d/100\) for each \(i\ge\allowbreak 2\).
Positive classes are \(i\equiv1,\allowbreak 2\pmod4\), and
\fitdisplay{
\sum_{i\ge1}(D_i)_+\le\frac{\delta_3+\delta_4}{1-\theta^2}
\le\frac{\theta^2}{1-\theta}d=169d/608<7d/25.
}{}
These bound the finite Duhamel sums. With
\(\sigma(\bar f_a)\le\allowbreak 6\sqrt d/5\) and \(\|R\|\le\allowbreak \sqrt d/5\),
\eqref{eq:Q95:2.1} gives \(\delta_+\le\allowbreak 72d/125\).
Use \eqref{eq:Q95:2.2} with \(n=\allowbreak a+L\) and the same \(V,\allowbreak N,\allowbreak c_L\)
to apply \eqref{eq:Q95:2.3}.
For \(L\ge\allowbreak 2\), \(\|V\|\le\allowbreak (113/100)\sqrt d\), and
\fitdisplay{
c_L+\E N\le(144/125+14/25+7/100-19/16)d=1189d/2000.
}{}
The total cost is \(264633d/10000\), whereas
\(8\Gamma=\allowbreak 276000d/10000\), leaving loss \(11367d/10000>\allowbreak d\).
For \(L=\allowbreak 1\), \(\beta_{a+1}+\beta_{a+2}\ge\allowbreak 51d/32\),
\(\|V\|\le\allowbreak (7/5)\sqrt d\), and
\(c_1+\E N\le\allowbreak (72/125+1/6-51/32)d\).
The cost is \(39767d/1500\), while \(8\Gamma=\allowbreak 41400d/1500\), leaving
\(1633d/1500>\allowbreak d\). Thus \(\E H^8\le\allowbreak  e^{-d}\) for all these triples.

\paragraph{Q95.4 The triple \((2,\allowbreak 3,\allowbreak 4)\).}
Use \(F=\allowbreak \bar f_2f_3h_4\). Its finite expansion is
\fitdisplay{
f_3+h_4=2-3d+dg-K^4\one,\qquad
R=dg-K^4\one,\qquad F\le\bar f_2e^{R-3d}.
}{}

\paragraph{Q95.4.1 The domain \(1/6\le\allowbreak  d\le\allowbreak 1/3\).}
The row and spectral bounds give \(|K^4\one|\le\allowbreak \theta\sqrt d\le\allowbreak  d\)
because \(\theta^2<\allowbreak 1/6\). Thus \(R\le\allowbreak 2d\), \(f_3h_4\le\allowbreak  e^{-d}\),
and \(H\le\allowbreak  e^{-d}p\bar f_{L+2}\).
Finite convolution and \eqref{eq:Q95:1.3} give
\(\E[(p\bar f_n)^8]\le\allowbreak  e^{253d/32}\) for \(n\ge\allowbreak 3\):
at \(n=\allowbreak 3\), the exponent is at most
\(-16d+17(d+\delta_2)\le\allowbreak 253d/32\), and at \(n\ge\allowbreak 4\), it is at most
\(-16\kappa d+17d=\allowbreak 15d/2\le\allowbreak 253d/32\).
Consequently \(\E H^8\le\allowbreak  e^{-3d/32}\) throughout this high domain.

\paragraph{Q95.4.2 The first step for \(0<\allowbreak d\le\allowbreak 1/6\).}
Put \(\mu=\allowbreak \E\bar f_2=\allowbreak 1-d\). Parity gives
\(\E R=\allowbreak -\delta_2\) and \(\E R^2=\allowbreak \delta_4+d^3\).
The pointwise range \(|R|\le\allowbreak  d+\theta\sqrt d<\allowbreak 1/3\) permits the same
Taylor estimate. For \(t=\allowbreak \delta_2/d\in[d,\allowbreak \theta]\),
\(\sigma(\bar f_2)^2=\allowbreak (1+t-d)d\) and
\(\E R^2\le\allowbreak (\theta^2t+d^2)d\), so
\fitdisplay{
\delta/d\le-(1-d)t+2\sqrt{1+t-d}\sqrt{\theta^2t+d^2}
+(12/5)(\theta^2t+d^2).
}{}
Apply \(2uv\le\allowbreak  u^2/3+3v^2\). The resulting bound is
\fitdisplay{
1/3-d/3+(27/5)d^2+t[-2/3+d+(27/5)\theta^2].
}{}
The coefficient of \(t\) is positive, its value at \(d=\allowbreak 0\) being
\(3449/15360\). At \(t=\allowbreak \theta\), the expression increases in \(d\)
because \(\theta-1/3=\allowbreak 7/96>\allowbreak 0\). At \(d=\allowbreak 1/6\), it is
\(1/16+59319/163840+7/576+3/20<\allowbreak 3/5\), with positive margin
\(19601/1474560\). Hence \(\delta_+\le\allowbreak 3d/5\).

\paragraph{Q95.4.3 The signed propagation means.}
For \(D_i=\allowbreak \E K'^iR\), the first values are
\fitdisplay{
\begin{gathered}
D_0=-\delta_2,\quad D_1=d^2,\quad D_2=\delta_3,\quad
D_3=-d\delta_2,\quad D_4=-\delta_4,\\
D_5=d\delta_3,\qquad D_6=\delta_5.
\end{gathered}
}{}
Positive contributions continue in precisely two classes:
\(i=\allowbreak 1+4j\) contributes \(d\delta_{1+2j}\), and \(i=\allowbreak 2+4j\) contributes
\(\delta_{3+2j}\). Therefore
\fitdisplay{
\sum_{i\ge1}(D_i)_+\le\frac{d^2+\delta_3}{1-\theta^2}
\le1019d/2565<2d/5.
}{}
For \(i\ge\allowbreak 5\), positive terms are at most
\(\max(d\delta_3,\allowbreak \delta_5)<\allowbreak d/36\), since \(d\le\allowbreak 1/6\) and
\(\theta^2<\allowbreak 1/6\). Also \(D_2<\allowbreak d/6\), while \(D_3,\allowbreak D_4\le\allowbreak 0\).
The positive partial sum used for \(L=\allowbreak 2\) is \(d^2\le\allowbreak  d/6\), and
for \(L=\allowbreak 3,\allowbreak 4\) it is \(d^2+\delta_3\le\allowbreak  d/3\).

\paragraph{Q95.4.4 Finite predecessors for every \(L\ge\allowbreak 2\).}
Here \(n=\allowbreak L+2\), and
\fitdisplay{
V=Q_L(K')g',\qquad Q_L(z)=G_{L+1}(z)-\mu d z^{L-1}-\mu z^{L+2}.
}{}
Write \(x=\allowbreak -z^2\in[0,\allowbreak \theta]\), \(\gamma=\allowbreak \mu d\in[0,\allowbreak 5/36]\),
and \(\mu\in[5/6,\allowbreak 1]\). The real coordinates below are nonnegative.
For \(L=\allowbreak 2\),
\(\Re Q=\allowbreak 1-x-\mu x^2\in[0,\allowbreak 1-x]\) and
\(\Im Q=\allowbreak (1-\gamma)\sqrt x\); hence
\(|Q|^2\le\allowbreak (1-x)^2+x\le\allowbreak 1\).
For \(L=\allowbreak 3\),
\fitdisplay{
\begin{gathered}
\Re Q=1-(1-\gamma)x\le1-31x/36,\\
\Im Q=\sqrt x(1-x-\mu x^2)\in[0,\sqrt x(1-x)],\\
|Q|^2-1\le-13x/18-1631x^2/1296+x^3\le0\quad(x\le1/2).
\end{gathered}
}{}
For \(L=\allowbreak 4\), put \(R_*=\allowbreak 1-x+x^2+x^3\). Then
\(\Re Q=\allowbreak 1-x+x^2+\mu x^3\le\allowbreak  R_*\) and
\(\Im Q\le\allowbreak \sqrt x(1-31x/36)\).
Since \(R_*^2=\allowbreak 1-2x+3x^2-x^4+2x^5+x^6\),
\fitdisplay{
|Q|^2-1\le-x+(23/18)x^2+(961/1296)x^3-x^4+2x^5+x^6.
}{}
Drop \(-x^4\) and divide by \(x\). On \(x\le\allowbreak 1/2\), the positive
terms sum to at most \(23/36+961/5184+1/8+1/32=\allowbreak 5083/5184<\allowbreak 1\).
For \(L=\allowbreak 5\),
\fitdisplay{
\begin{gathered}
\Re Q=1-x+(1-\gamma)x^2\le1-x+x^2,\\
\Im Q=\sqrt x(1-x+x^2+\mu x^3)\le\sqrt xR_*,\\
|Q|^2-1\le-x+x^2+x^3+x^4-x^5+2x^6+x^7.
\end{gathered}
}{}
After \(-x^5\) is dropped and division by \(x\), the positive terms
are at most \(1/2+1/4+1/8+1/16+1/64=\allowbreak 61/64<\allowbreak 1\).
For every \(L\ge\allowbreak 6\), finite geometric summation and
\((1+x)^{-1/2}\le\allowbreak 1-x/2+3x^2/8\) give
\fitdisplay{
|Q_L|\le1-x/2+3x^2/8+x^{7/2}+(5/36)x^{5/2}+x^4.
}{}
Divide the positive error by \(x\) and use \(x\le\allowbreak 5/12\),
\(\sqrt x\le\allowbreak 2/3\). Its four terms are at most
\(5/32+25/216+25/648+125/1728=\allowbreak 1985/5184<\allowbreak 2/5\).
Thus \(|Q_L|\le\allowbreak 1-x/10\le\allowbreak 1\). At \(x=\allowbreak 0\), equality holds.
For \(L=\allowbreak 1\), separately, \(Q_1=\allowbreak 1-\gamma+z-\mu z^3\), and
\(|Q_1|^2\le\allowbreak 1+x(1+x)^2\le\allowbreak 1+26325/32768<\allowbreak 29/16\).

\paragraph{Q95.4.5 The low-domain moment bounds.}
Use \(\Gamma=\allowbreak 3d\), the constants in Q95.2 and \(\|V\|^2\le\allowbreak  d\)
for \(L\ge\allowbreak 2\). For \(L=\allowbreak 2\),
\(c_L+\E N\le\allowbreak (6/5+1/3+1/6-19/16)d=\allowbreak 41d/80\), giving exponent
at most \(-29d/10\). For \(L=\allowbreak 3,\allowbreak 4\), the positive partial sum is
at most \(d/3\) and \(D_L\le\allowbreak 0\), so
\(c_L+\E N\le\allowbreak (6/5+2/3-19/16)d=\allowbreak 163d/240\), giving \(-47d/30\).
For every \(L\ge\allowbreak 5\),
\(c_L+\E N\le\allowbreak (6/5+4/5+1/36-19/16)d=\allowbreak 121d/144\), giving \(-5d/18\).
For \(L=\allowbreak 1\), \(\beta_3+\beta_4=\allowbreak 2d\), \(D_1=\allowbreak d^2\),
\(c_1+\E N\le\allowbreak -37d/30\), and the variance cost is at most \(493d/16\).
The resulting exponent is at most \(-733d/240\).
The low domain therefore gives \(\E H^8\le\allowbreak  e^{-5d/18}\).
It meets the high domain at \(d=\allowbreak 1/6\) with the same
\(F=\allowbreak \bar f_2f_3h_4\). The zero boundary and positive restoration
of decorations give \(\E H^8\le\allowbreak  e^{-3d/32}\) for every \(L\ge\allowbreak 1\).

\paragraph{Q95.5.1 Finite tails for the first short pair three.}
Let \(2<\allowbreak 3<\allowbreak c_1<\allowbreak b_2<\allowbreak c_2<\allowbreak \cdots<\allowbreak b_k<\allowbreak c_k\), \(k\ge\allowbreak 2\), and
\fitdisplay{
F=\bar f_2f_3h_{c_1}\prod_{i=2}^kh_{b_i}f_{c_i},\qquad
S=\sum_{i=2}^k\sum_{r=b_i+1}^{c_i}K^r\one.
}{}
Every later \(b_i\ge\allowbreak 5\), so its pair has loss \(16d/5\).
The finite zero-one mask \(S\) has first possible exponent \(c_1+2\),
and \(\|S\|^2\le\allowbreak \theta^{b_2}C_*d\),
\(|S|^2\le\allowbreak \theta^{b_2-1}C_*d\), where
\(C_*=\allowbreak 32768/16245<\allowbreak 100/49\) is \eqref{eq:L90:masknorm}.
Since \(\theta^3<\allowbreak 7/100\), \(\sqrt\theta<\allowbreak 13/20\), and
\(\theta^5C_*<\allowbreak 9/400\), we have
\fitdisplay{
\begin{array}{c|c|c|c}
c_1&b_2\text{ lower}&\|S\|/\sqrt d\text{ upper}&|S|/\sqrt d\text{ upper}\\\hline
4&5&3/20&1/4\\
5&6&1/10&3/20\\
\ge6&\ge7&7/100&1/10
\end{array}
}{}
For instance \(\theta^6C_*<\allowbreak (7/100)^2(100/49)=\allowbreak 1/100\), and the next
norm is below \(13/200<\allowbreak 7/100\); pointwise bounds shift the power by one.
In all cases \(\E S\le\allowbreak \delta_4/(1-\theta^2)<\allowbreak d/12\) and
\(D_1(S)<\allowbreak d/30\). The complete positive \(D_i(S)\) tails are bounded
respectively by \(d/5,\allowbreak d/8,\allowbreak d/12\). For exponents \(r\ge\allowbreak 6\), the first
four positive moment indices are \(\delta_5,\allowbreak \delta_4,\allowbreak \delta_5,\allowbreak \delta_6\),
giving \(\theta^3d/(1-\theta)^2=\allowbreak 2197d/11552<\allowbreak d/5\).
For \(r\ge\allowbreak 7\), they are \(\delta_5,\allowbreak \delta_6,\allowbreak \delta_5,\allowbreak \delta_6\),
giving \(2\theta^4C_*d<\allowbreak (1/18)(100/49)d=\allowbreak 50d/441<\allowbreak d/8\).
For \(r\ge\allowbreak 8\), they are \(\delta_5,\allowbreak \delta_6,\allowbreak \delta_7,\allowbreak \delta_6\),
giving \(\theta^4d/(1-\theta)^2=\allowbreak 28561d/369664<\allowbreak d/12\).
For each \(i\ge\allowbreak 2\), the individual bounds are \(d/12,\allowbreak d/30,\allowbreak d/72\),
respectively. The middle bound follows from
\(\delta_5/(1-\theta^2)<\allowbreak (1/36)/(5/6)d=\allowbreak d/30\);
the last from \(\delta_6/(1-\theta^2)<\allowbreak (1/90)/(5/6)d=\allowbreak d/75<\allowbreak d/72\).

\paragraph{Q95.5.2 The first long arm \(c_1=\allowbreak 4\).}
Take \(R_0=\allowbreak dg-K^4\one\), \(R=\allowbreak R_0+S\), and
\(\Gamma=\allowbreak [3+(16/5)(k-1)]d\). First let \(d\le\allowbreak 1/6\).
The pointwise range is less than
\(1/3+(1/4)\sqrt{1/6}<\allowbreak 1/2\), so the same Taylor bound applies.
Adding \(S\) increases the first-step constant by at most
\fitdisplay{
(\E S)_++2\sigma(\bar f_2)\|S\|
+(12/5)(2\|R_0\|\|S\|+\|S\|^2).
}{}
Since \(\|R_0\|^2<\allowbreak d/10<\allowbreak d/9\) and
\(\sigma(\bar f_2)\le\allowbreak 6\sqrt d/5\),
\fitdisplay{
\delta_+/d\le3/5+1/12+9/25+(12/5)(1/10+9/400)
=1003/750<27/20.
}{}
The positive \(D_i\) tail is at most \(3d/5\), and each
\(D_L\le\allowbreak  d/4\) for \(L\ge\allowbreak 2\).
The base \(Q_L\) of Q95.4.4 is a contraction for \(L\ge\allowbreak 2\); adding
the propagated \(S\) gives \(\|V\|\le\allowbreak (11/10)\sqrt d\).
For \(k=\allowbreak 2,\allowbreak L\ge\allowbreak 2\),
\(c_L+\E N\le\allowbreak (27/10+6/5-19/16+1/4)d=\allowbreak 237d/80\).
The cost apart from \(-8\Gamma\) is \(4427d/100\), whereas
\(8\Gamma=\allowbreak 4960d/100\); the net loss is \(533d/100>\allowbreak 5d\).
For \(L=\allowbreak 1\), \(\|Q_1g'\|\le\allowbreak (27/20)\sqrt d\); adding \(S\) gives
\(\|V\|\le\allowbreak 3\sqrt d/2\). Also
\(D_1\le\allowbreak  d^2+d/30\le\allowbreak  d/5\), \(\beta_3+\beta_4=\allowbreak 2d\), and
\(c_1+\E N\le\allowbreak (27/20+1/5-2)d=\allowbreak -9d/20\).
The net loss for \(k=\allowbreak 2\) is \(299d/20>\allowbreak 5d\).
For \(d\ge\allowbreak 1/6\), the first pair is at most \(e^{-d}\), while
\(|S|\le\allowbreak (1/4)\sqrt d\le\allowbreak 5d/8\). The pair product is at most
\(e^{-[3/8+(16/5)(k-1)]d}\).
The bare-path cost \(253d/32\) leaves loss \(3311d/160>\allowbreak 5d\) at \(k=\allowbreak 2\).
In all cases an additional pair changes only \(\Gamma\). Therefore
\fitdisplay{
\E H^8\le e^{-[5+(128/5)(k-2)]d}\qquad(c_1=4,\ k\ge2,\ L\ge1).
}{}

\paragraph{Q95.5.3 The exact base for \(c=\allowbreak c_1\ge\allowbreak 5\) and \(d\le\allowbreak 1/32\).}
Set \(m=\allowbreak c-3\ge\allowbreak 2\), \(B=\allowbreak \beta_{c-1}\), \(\mu=\allowbreak 1-d\), and
\fitdisplay{
W=\sum_{r=4}^cK^r\one=K^3G_m(K)g,\qquad
R_0=Bg-W,\qquad\Gamma_0=d+\beta_c+B.
}{}
The finite difference gives \(f_3+h_c\le\allowbreak 2-\Gamma_0+R_0\), so
\(F\le\allowbreak \bar f_2e^{R_0+S-\Gamma_0-(16/5)(k-1)d}\).
For \(c=\allowbreak 5\), \(B=\allowbreak d\) and \(\Gamma_0=\allowbreak 3d-\delta_2\ge\allowbreak 83d/32\).
For \(c\ge\allowbreak 6\), \(\Gamma_0\ge\allowbreak  d+2\tau\ge\allowbreak 35d/16\).
Consecutive alternating moments give \(\E W\ge\allowbreak 0\) and \(\E(gW)\ge\allowbreak 0\).
More precisely \(\E W=\allowbreak \delta_2\) for \(c=\allowbreak 5\), while
\(\E W\ge\allowbreak \delta_2-\delta_3\ge\allowbreak \kappa\delta_2\) for \(c\ge\allowbreak 6\).
Thus \(\E R_0=\allowbreak -\delta_2\) or at most \(-\kappa\delta_2\), respectively,
and \(\E R_0^2\le\allowbreak  B^2d+\E W^2\). For \(t=\allowbreak \delta_2/d\),
\fitdisplay{
\E R_0^2/d\le
\begin{cases}\theta^2(1+\theta)t+d^2,&c=5,\\
\theta^2t+d^2,&c\ge6.
\end{cases}
}{}
For \(c=\allowbreak 5\), \(\E R_0^2<\allowbreak d/10\), since
\(\theta^3(1+\theta)+1/1024=\allowbreak 99889/1048576<\allowbreak 1/10\).
Also \(\|W\|^2=\allowbreak \delta_4+\delta_5\le\allowbreak 98865d/1048576<\allowbreak 25d/256\),
so \(\|W\|<\allowbreak 5\sqrt d/16\).
For \(c\ge\allowbreak 6\), the variance is below \(9d/100\), using
\(\theta^3<\allowbreak 7/100\) and \(1/1024<\allowbreak 1/100\).
Including \(S\), the pointwise range is at most
\(d+(51/80)\sqrt d<\allowbreak 1/32+51/400=\allowbreak 127/800<\allowbreak 1/3\).

Both base first-step constants have \((\delta_0)_+\le\allowbreak 3d/5\).
For \(c\ge\allowbreak 6\), apply \(2uv\le\allowbreak  u^2/3+3v^2\). The bound divided by
\(d\) is
\fitdisplay{
1/3-d/3+(27/5)d^2+t[1/3-(1-d)\kappa+(27/5)\theta^2].
}{}
Its \(t\)-coefficient is positive. At \(t=\allowbreak \theta\), the linear
\(d\)-coefficient \(\kappa\theta-1/3\) is negative; dropping it and
using \(d^2\le\allowbreak 1/1024\) yields \(292389/491520<\allowbreak 3/5\), with margin
\(2523/491520\).
For \(c=\allowbreak 5\), the corresponding expression is
\fitdisplay{
1/3-d/3+(27/5)d^2+t[-2/3+d+(27/5)\theta^2(1+\theta)].
}{}
It increases in \(t\). At \(t=\allowbreak \theta\), the linear coefficient is
\(7/96>\allowbreak 0\); at \(d=\allowbreak 1/32\), the bound is
\(9109889/15728640<\allowbreak 3/5\), with margin \(327295/15728640\).
For the base signed means, \(-W\) contributes positive classes
\(i\equiv2,\allowbreak 3\pmod4\), totaling at most
\((\delta_3+\delta_4)/(1-\theta^2)\).
The \(Bg\) term contributes class \(i\equiv1\pmod4\), at most
\(Bd/(1-\theta^2)\). Hence
\fitdisplay{
\sum_{i\ge1}(D_i(R_0))_+\le169d/608+32d/855<d/3.
}{}
Also \(D_1(R_0)\le\allowbreak  Bd\le\allowbreak  d^2\), and every
\(D_i(R_0)\le\allowbreak \delta_3<\allowbreak d/6\) for \(i\ge\allowbreak 2\).
At \(i=\allowbreak 3\), the \(Bg\) term is negative; at \(i=\allowbreak 5\), dropping the
negative \(W\) term leaves \(B\delta_3\le\allowbreak \delta_3\).

\paragraph{Q95.5.4 A finite predecessor for all \(c\ge\allowbreak 5,\allowbreak L\ge\allowbreak 2\) on the small domain.}
Put \(z=\allowbreak K'\) and \(\gamma=\allowbreak \mu B\le\allowbreak 1/32\). The base predecessor is
\fitdisplay{
Q_{L,c}(z)=G_{L+1}(z)-\gamma z^{L-1}-\mu z^{L+2}G_m(-z).
}{}
Write \(G_m(-z)=\allowbreak U-zV\). For \(m\ge\allowbreak 2\), both \(U,\allowbreak V\) are finite
alternating sums in \([1-x,\allowbreak 1]\); also \(x\le\allowbreak \theta<\allowbreak 1/2\) and
\(\mu\ge\allowbreak 31/32\).
For \(L=\allowbreak 2\), the real coordinate is positive and
\(1-x-\mu x^2U\le\allowbreak 1-x-(31/64)x^2\).
The imaginary coordinate \(\sqrt x[1-\gamma+\mu x^2V]\) is positive
and at most \(\sqrt x(1+x^2)\). Thus
\fitdisplay{
|Q|^2-1\le-x+(1/32)x^2+(95/32)x^3+(1/4)x^4+x^5.
}{}
After division by \(x\), the positive terms sum to at most
\(109/128<\allowbreak 1\) on \(x\le\allowbreak 1/2\).
For \(L=\allowbreak 3\), \(|G_4|\le\allowbreak (1-x)(1+x/2)=\allowbreak 1-x/2-x^2/2\).
Adding \(\gamma x\) and \((6/5)x^{5/2}\) gives
\(|Q|\le\allowbreak 1-15x/32+3x^2/10\le\allowbreak 1-51x/160\le\allowbreak 1\).
For \(L=\allowbreak 4\), \(|G_5|\le\allowbreak 1-x/2+x^2/2\), since the squared
difference is \(x^2(1+2x-3x^2)/4\ge\allowbreak 0\). Adding
\(\gamma x^{3/2}\) and \((6/5)x^3\) gives
\(|Q|\le\allowbreak 1-23x/48+x^2/2+(6/5)x^3\le\allowbreak 1-x/16\), using
\(x/2+(6/5)x^2\le\allowbreak 5/12\) for \(x\le\allowbreak 5/12\).
For every \(L\ge\allowbreak 5\), finite geometric summation and the negative
third derivative of \((1+x)^{-1/2}\) give
\fitdisplay{
|G_{L+1}|\le\frac{1+x^3}{\sqrt{1+x}}\le1-x/2+3x^2/8+x^3.
}{}
The extra terms are at most \(x^2/32+(6/5)x^{7/2}\).
After division by \(x\), use \(x\le\allowbreak 5/12\), \(\sqrt x\le\allowbreak 2/3\) to bound
\((13/32)x+x^2+(6/5)x^{5/2}\le\allowbreak 185/384<\allowbreak 1/2\).
Therefore \(|Q|\le\allowbreak 1-7x/384\le\allowbreak 1\); equality at \(x=\allowbreak 0\) is direct.
For \(L=\allowbreak 1\), the safe norm bounds are \(31/20\) for \(c=\allowbreak 5\) and
\(3/2\) for \(c\ge\allowbreak 6\). These use \(\|G_2\|\le\allowbreak 6/5\),
\(\gamma\le\allowbreak 1/32\), and \(\|W\|/\sqrt d<\allowbreak 5/16\) or
\(\theta^{3/2}<\allowbreak (13/32)(13/20)\), respectively.

\paragraph{Q95.5.5 Small-domain moment bounds.}
For \(c=\allowbreak 5\), \(\|R_0\|<\allowbreak \sqrt d/3\), \(\|S\|\le\allowbreak \sqrt d/10\).
The increase of \(\delta/d\) is at most
\(1/12+6/25+(12/5)(1/15+1/100)=\allowbreak 761/1500\), giving
\(\delta_+\le\allowbreak 9d/8\).
The positive \(D_i\) tail is at most \(d/3+d/8=\allowbreak 11d/24\), and
\(D_L\le\allowbreak  d/6+d/30=\allowbreak d/5\) for \(L\ge\allowbreak 2\).
The predecessor bound is \(\|V\|\le\allowbreak (16/15)\sqrt d\), so
\(c_L+\E N\le\allowbreak 523d/240\).
The cost is \(523d/30+4352d/225\); for \(k=\allowbreak 2\),
\(8\Gamma\ge\allowbreak 927d/20\), leaving \(17234d/1800>\allowbreak 9d\).
At \(L=\allowbreak 1\), \(\|V\|\le\allowbreak 33\sqrt d/20\) and
\(c_1+\E N\le\allowbreak (9/8+1/32+1/30-2)d=\allowbreak -389d/480\);
the net loss is \(7861d/1200>\allowbreak 6d\).

For \(c\ge\allowbreak 6\), \(\|R_0\|<\allowbreak 3\sqrt d/10\), \(\|S\|\le\allowbreak 7\sqrt d/100\),
and
\fitdisplay{
\delta\le[3/5+1/12+21/125+(12/5)(21/500+49/10000)]d<d.
}{}
The positive tail is at most \(d/3+d/12=\allowbreak 5d/12\), while
\(D_L\le\allowbreak  d/6+d/72<\allowbreak d/5\) for \(L\ge\allowbreak 2\).
Now \(\|V\|\le\allowbreak 21\sqrt d/20\) and \(c_L+\E N\le\allowbreak 443d/240\), with
total cost \(443d/30+7497d/400\). Since \(8\Gamma\ge\allowbreak 431d/10\),
the net loss is \(11509d/1200>\allowbreak 9d\).
At \(L=\allowbreak 1\), \(\|V\|\le\allowbreak 8\sqrt d/5\) and
\(c_1+\E N\le\allowbreak (1+1/32+1/30-2)d=\allowbreak -449d/480\), leaving
\(2119d/300>\allowbreak 7d\). Additional pairs increase only \(8\Gamma\), by
\(128d/5\) for each pair.

\paragraph{Q95.5.6 The domain \(d\ge\allowbreak 1/32\).}
Here \(\sqrt d\le\allowbreak \sqrt{32}\,\allowbreak d<\allowbreak (23/4)d\).
For \(c=\allowbreak 5\),
\fitdisplay{
|R_0+S|\le d+[(6/5)\theta+3/20]\sqrt d
\le d+(51/80)(23/4)d=1493d/320.
}{}
At \(k=\allowbreak 2\), \(\Gamma\ge\allowbreak 927d/160\), so the pair product is at most
\(e^{-361d/320}\). Multiplying the bare path \(p\bar f_{L+2}\) and
using its \(253d/32\) cost leaves \(179d/160>\allowbreak d\), for all \(L\ge\allowbreak 1\).
For \(c\ge\allowbreak 6\), use \(\|G_m\|\le\allowbreak 1\), obtaining
\(|R_0+S|\le\allowbreak  d+(\theta+1/10)\sqrt d\).
Even \(\sqrt d<\allowbreak 6d\) gives \(|R_0+S|\le\allowbreak 323d/80\).
At \(k=\allowbreak 2\), \(\Gamma\ge\allowbreak 431d/80\), so the pair product is at most
\(e^{-27d/20}\); the bare-path cost leaves \(463d/160>\allowbreak 2d\).
Both domains include \(d=\allowbreak 1/32\) and keep the same direction. Thus
\fitdisplay{
\boxed{\E H^8\le e^{-[1+(128/5)(k-2)]d}
\qquad(c_1\ge5,\ k\ge2,\ L\ge1).}
}{}
Together with Q95.5.2, the weaker bound holds for every \(c_1\ge\allowbreak 4\).

\paragraph{Q95.6 Relation to the ordinary rows.}
Appendix~\ref{proof:Q5-C093} proves the direction
\(\bar f_a\prod_i h_{b_i}f_{c_i}\), with bounds
\(e^{-[10+(128/5)(k-3)]d}\) for \(a=\allowbreak 2,\allowbreak b_1=\allowbreak 3,\allowbreak k\ge\allowbreak 3\),
\(e^{-[17+(128/5)(k-2)]d}\) for \(a=\allowbreak 2,\allowbreak 3,\allowbreak b_1=\allowbreak 4,\allowbreak k\ge\allowbreak 2\), and
\(e^{-[1+(128/5)(k-1)]d}\) for \(a\ge\allowbreak 2,\allowbreak b_1\ge\allowbreak 5,\allowbreak k\ge\allowbreak 1\).
Q95.3 supplies \(b_1=\allowbreak 4,\allowbreak k=\allowbreak 1\), Q95.4 supplies \((2,\allowbreak 3,\allowbreak 4)\), and
Q95.5 supplies \(b_1=\allowbreak 3,\allowbreak k\ge\allowbreak 2\) with first pair \(f_3h_{c_1}\).
These integer ranges leave only the triple \((2,\allowbreak 3,\allowbreak c)\), \(c\ge\allowbreak 5\),
for the calculation of Q5-C097.
For a list belonging to more than one row, choose one row from its
integer structure and retain that row's full direction and moment bound.
Positive comparison restores every genuine equal pair at its vertex.
Replacing \(K\) by \(-K\) reverses all arms, connector, receiving
unit arm and decorations together, with the fixed measure and
all \(\delta_i\) preserved.
